%! Semi-log Plots — Unit 5, Lesson 5.8 — Homework
\documentclass[10pt]{article}
\usepackage{precalculus-article}
\usepackage{precalculus-boxes}
\usepackage{grading}
\gradingassignment{precalc-u05-l08-hw}
\gradingcourse{precalculus}
\newcommand{\TallMath}[1]{$\displaystyle #1\rule[-1.4em]{0pt}{3.2em}$}

\begin{document}

\pageheader{Unit 5, Lesson 5.8}{Homework \quad\textperiodcentered\quad Write each final answer in its box. Only what is inside a box is graded.}

\vspace{0.1in}

\begin{enumerate}[leftmargin=1.8em, itemsep=14pt]

\item \textbf{Conceptual.}  Explain in your own words why exponential data appears linear
      on a semi-log plot.  Your explanation should reference what happens algebraically when
      you apply a logarithm to both sides of $y = ab^x$.  (No equations required; write
      in complete sentences.)

      \rubricbox{\gradedblank[id=hw1, type=rubric, points=2,
  display={Taking $\log$ of both sides of $y = ab^x$ brings the exponent down: $\log y = x\log b + \log a$, a linear function of $x$. So $\log y$ vs.\ $x$ is a straight line; multiplicative growth in $y$ becomes additive growth in $\log y$.}]{\linewidth}{42pt}}

\item Given the exponential model $y = 2 \cdot 5^x$, write the linearized form using
      $\log_{10}$.  Identify the slope and $y$-intercept of the linearized model.

\vspace{4pt}
Linearized form: $\log_{10}(y) = $ \smash{\gradedblank[id=hw2a, type=short, points=1, answer={(\log_{10} 5) x + \log_{10} 2}, display={$(\log_{10} 5)\cdot x + \log_{10} 2 \approx 0.699x + 0.301$}]{6cm}{13pt}}

\vspace{4pt}
Slope $= $ \smash{\gradedblank[id=hw2b, type=short, points=1, answer={\log_{10} 5}, display={$\log_{10}(5) \approx 0.699$}]{4cm}{13pt}} \qquad $y$-intercept $= $ \smash{\gradedblank[id=hw2c, type=short, points=1, answer={\log_{10} 2}, display={$\log_{10}(2) \approx 0.301$}]{4cm}{13pt}}

\item A semi-log plot of $(x,\,\log_{10}(y))$ data shows a linear pattern with slope $0.477$
      and $y$-intercept $1.0$.  Write the exponential model $y = a\cdot b^x$.

\vspace{4pt}
$b = 10^{0.477} \approx $ \smash{\gradedblank[id=hw3a, type=short, points=1, answer={3}]{2cm}{13pt}} \qquad $a = 10^{1.0} = $ \smash{\gradedblank[id=hw3b, type=short, points=1, answer={10}]{2cm}{13pt}}

\vspace{4pt}
Exponential model: $y = $ \smash{\gradedblank[id=hw3c, type=short, points=1, answer={10 \cdot 3^x}]{5cm}{13pt}}

\item The table below gives data for a population $P$ over time $t$.

\vspace{4pt}
\begin{center}
\renewcommand{\arraystretch}{2.0}
\begin{tabular}{|c|c|c|c|c|}
\hline
\rowcolor{sky}
$t$ & $0$ & $2$ & $4$ & $6$ \\
\hline
$P$ & $500$ & $2000$ & $8000$ & $32000$ \\
\hline
$\log_{10}(P)$ & \smash{\gradedblank[id=hw4a1, type=short, check=numeric, tolerance=0.0005, points=1, answer={2.699}]{1.8cm}{15pt}} & \smash{\gradedblank[id=hw4a2, type=short, check=numeric, tolerance=0.0005, points=1, answer={3.301}]{1.8cm}{15pt}} & \smash{\gradedblank[id=hw4a3, type=short, check=numeric, tolerance=0.0005, points=1, answer={3.903}]{1.8cm}{15pt}} & \smash{\gradedblank[id=hw4a4, type=short, check=numeric, tolerance=0.0005, points=1, answer={4.505}]{1.8cm}{15pt}} \\
\hline
\end{tabular}
\end{center}

\begin{enumerate}[leftmargin=1.5em, itemsep=6pt, label=\alph*.]
  \item Complete the $\log_{10}(P)$ row (round to three decimal places).
  \item Find the slope of $\log_{10}(P)$ vs.\ $t$ using any two consecutive points.
        \quad Slope $= $ \smash{\gradedblank[id=hw4b, type=short, check=numeric, tolerance=0.0005, points=1, answer={0.301}]{4cm}{13pt}}
  \item Find $b = 10^{\text{slope}} \approx $ \smash{\gradedblank[id=hw4c1, type=short, points=1, answer={2}]{2.5cm}{13pt}} and $a = 10^{\log_{10}(P_0)} \approx $ \smash{\gradedblank[id=hw4c2, type=short, points=1, answer={500}]{2.5cm}{13pt}}.
  \item Write the exponential model: $P(t) = $ \smash{\gradedblank[id=hw4d, type=short, points=1, answer={500 \cdot 2^t}]{5cm}{13pt}}
  \item Verify: $P(2) = $ \smash{\gradedblank[id=hw4e1, type=short, points=1, answer={2000}]{4cm}{13pt}}. Does this match the table? \smash{\gradedblank[id=hw4e2, type=short, check=text, points=1, answer={yes}]{1.5cm}{13pt}}
\end{enumerate}

\item \textbf{Comparison.}  Two scientists analyze the same data set.  Scientist A adds a
      constant to all $y$-values before plotting on a regular graph to check for exponential
      patterns.  Scientist B uses a semi-log plot instead.  What advantage does Scientist B's
      approach have?  Reference EK 2.15.A.2 in your response.

      \rubricbox{\gradedblank[id=hw5, type=rubric, points=2,
  display={B needs no guessed constant: by EK 2.15.A.2 a constant never has to be added to the $y$-values on a semi-log plot to show an exponential model fits. If $\log y$ vs.\ $x$ is linear, the model is exponential.}]{\linewidth}{42pt}}

\item \textbf{AP-style multiple choice.} Write the letter of the single best answer in its Letter box.  For the exponential model $y = 8 \cdot 6^x$,
      which of the following gives the slope of the linearized semi-log model using
      base-10 logarithms?

\vspace{6pt}
\begin{enumerate}[leftmargin=2em, itemsep=4pt, label=(\Alph*)]
  \item $\log_{10}(8)$
  \item $\log_{10}(6)$
  \item $\log_{10}(48)$
  \item $6$ \hfill\textit{Letter:}~\gradedblank[id=hw6, type=choice, points=1, answer=B]{0.42in}{13pt}
\end{enumerate}

\end{enumerate}

\vspace{0.1in}

\begin{tcolorbox}[colback=goldbg, colframe=goldacc, arc=2mm, left=3mm, right=3mm,
    top=2mm, bottom=2mm, title=\textbf{Extension}, fonttitle=\bfseries]
Suppose you take the natural log (ln) of $y = ab^x$ instead of $\log_{10}$.  Write the
linearized model.  How do the slope and $y$-intercept differ from the base-10 version?
(Hint: use the change-of-base formula to express $\ln(b)$ in terms of $\log_{10}(b)$.)

\writelines{3}
\end{tcolorbox}

\end{document}
